\documentclass[12pt]{article}

\input{/home/iasbeck/preamble.tex}

\renewcommand{\headrulewidth}{0.5pt} % Redefine a espessura da lina no cabeçalho
\lhead{Minimização de uma função restrita pela aplicação do MPE} % Lado esquerdo do cabeçalho
\chead{} % Centro do cabeçalho
\rhead{Arthur Iasbeck} % Lado direito  do cabeçalho
\lfoot{} % Lado esquerdo do rodapé
\cfoot{} % Centro do rodapé
\rfoot{\thepage} % Lado direito  do rodapé

% Folha A4
% Margem superior = Margem esquerda = 3cm
% Margem direita = Margem inferior = 2cm
% Fonte Times New Roman 12
% Títulos de Figuras e Tabelas Times New Roman 10
% Espaçamento simples e alinhamento justificado

% Início do documento
\begin{document}
	\singlespacing % Espaçamento simples
	%\pagestyle{empty} % Remove numeração das páginas
	
	% Inserção da capa
	\thispagestyle{empty}
	\newgeometry{top=1.5cm, bottom=2cm, left=2cm, right=2cm} 
	
	\begin{figure}[H]
		\centering
		\includegraphics[width=\linewidth]{cabecalho.png}
	\end{figure}
	
	\begin{center}
		\large
		\textbf{UNIVERSIDADE FEDERAL DE UBERLÂNDIA} \\ \vspace{0.3cm}
		\textbf{FACULDADE DE ENGENHARIA MECÂNICA} \\~\\~\\~\\~\\~\\~\\~\\
		\Large
		\textbf{MINIMIZAÇÃO DE UMA FUNÇÃO RESTRITA A PARTIR DA APLICAÇÃO DO MÉTODO DA PENALIDADE EXTERNA} \\~\\~\\~\\~\\~\\~\\
		\large
		\textbf{ARTHUR HENRIQUE IASBECK} \\~\\~\\~\\~\\~\\~\\~\\
		\large
		\textbf{UBERLÂNDIA} \\
		\textbf{31 DE OUTUBRO DE 2019}
	\end{center}
	\newpage
	\restoregeometry
	
	\setcounter{page}{1} % Impede que a capa seja contada na paginação
	
	\section{OBJETIVO}
		Foi proposta no presenta trabalho a minimização da função $ f(x) $, Eq. \ref{fObj},
		\begin{equation}
			f(x) = \frac{1}{3}(x_1 + 1)^3 + x_2
			\label{fObj}
		\end{equation}
		
		sujeita às seguintes restrições
		\begin{gather}
			-x_1 + 1 \leq 0 \\
			-x_2 \leq 0
		\end{gather}
		
	\section{METODOLOGIA}
		Uma vez que o problema proposto contém restrições de desigualdade, é necessário tratá-las de forma a transformá-lo num problema irrestrito, de forma que seja possível empregar métodos como o das Variáveis Métricas, o de Newton ou o da Máxima Descida, por exemplo, para obtenção da solução. 
		
		No presente trabalho adotou-se o Método da Penalidade Externa (MPE) para tratamento das restrições. Este método se baseia na definição de uma nova função objetivo $ \phi(x) $, Eq. \ref{fObjPen}, que inclui uma parcela de penalização que garante que $ f(x) \rightarrow \infty $ caso $ x $ não atenda as restrições. 
		\begin{equation}
			\label{fObjPen}
			\phi(x) = f(x) + P(x)
		\end{equation}
		
		No MPE defini-se a função $ P(x) $ da seguinte forma
		\begin{equation}
			P(x) = r_p (\sum_{i = 1}^{m}(h_i(x))^2 + \sum_{j = 1}^{k} (max(0, g_j(x)))^2)
		\end{equation}		 
		
		\noindent em que $ m $ e $ k $ representam respectivamente o número de restrições de igualdade e desigualdade, e $ r_p $ é um valor que tende a infinito com o passar das iterações. Observa-se que $ P(x) $ será maior que zero caso $ h_i(x) \neq 0 $ ou $ g_j(x) > 0 $ (o que faria, neste último caso, com que $ max(0, g_j(x)) = g_j(x)$). Para evitar problemas de ordem numérica, é necessário que o valor de $ r_p $ seja inicialmente pequeno e cresça à medida que evolui o processo de otimização, de forma a garantir que $ \phi(x^*) = f(x^*) $ sendo $ x^* $ uma solução que atenda às restrições de igualdade e desigualdade.
		
		Aplicando a definição de $ P(x) $ ao problema proposto obtém-se
		\begin{equation}
			\phi(x) = \frac{1}{3}(x_1 + 1)^3 + x_2 + r_p(max(0, -x_1 + 1)^2 + max(0, -x_2)^2)
		\end{equation}  
		
		A determinação do mínimo de $ \phi(x) $ consiste num problema multi-dimensional irrestrito, ao qual foi aplicado o Método das Variáveis Métricas. Uma vez reduzida a ordem do problema, o mesmo pôde ser resolvido, em última instância, aplicando-se o Método da Seção Áurea. 
		
		A solução $ x^* $ foi obtida por meio do emprego de um código escrito no Matlab\textsuperscript{\textregistered}, que foi condensado na função
		
		\noindent \mcode{[xOpt, fOpt, nVal, k] = varMet(x0, theta, rp0, rpInc, tol, h)} \\
		
		\noindent em que 
		\begin{itemize}
			\item \mcode{xOpt}: Solução ótima obtida a partir da execução da função \mcode{varMet}.
			\item \mcode{fOpt}: Valor de $ f(x^*) $.
			\item \mcode{nVal}: Número de avaliações da função objetivo.
			\item \mcode{k}: Número de iterações.
			\item \mcode{x0}: Palpite inicial para a solução.
			\item \mcode{theta}: Valor de $ \theta $ a ser utilizado na computação do Método das Variáveis Métricas.
			\item \mcode{rp0}: Valor inicial de $ r_p $.
			\item \mcode{rpInc}: Incremento de $ r_p $, que é redefinido a cada iteração como \mcode{rp = rp*rpInc}.
			\item \mcode{tol}: Tolerância adotada para o critério de parada, que é executado até que a norma euclidiana entre $ x_{k+1} $ e $ x_k $ seja menor do que \mcode{tol}.
			\item \mcode{h}: Incremento adotado na computação numérica do gradiente da função objetivo.
 		\end{itemize}
 
 		Juntamente à função \mcode{varMet}, foram definidas outras três funções denominadas \mcode{funcFile.m}, \mcode{gradFile.m} e \mcode{funcFile.m}, nas quais o usuário deveria inserir, respectivamente, a função objetivo, o seu gradiente (caso o mesmo pudesse ser determinado analiticamente) e as restrições de igualdade e desigualdade atribuídas ao problema. Se o gradiente da função objetivo não for informado pelo usuário, será adotada a computação numérica do mesmo. Cabe ressaltar que as restrições devem ser informadas em sua forma matricial, como mostrado a seguir.
 		
 		\begin{lstlisting}
function [g, h] = constFile(x)

% Cada linha da matriz 'g' representa uma restricao de desigualdade
g = [-x(1) + 1
     -x(2)];

% Cada linha da matriz 'h' representa uma restricao de igualdade
h = [];
 		\end{lstlisting}
 
 		\vspace{0.4cm}
 		Executando a função \mcode{varMet} considerando as entradas apresentadas na Tab. \ref{entradasVarMet}, foi obtido o resultado apresentado abaixo, condizente com a solução analítica esperada.
 		
 		\begin{lstlisting}
x1* = 1.0000
x2* = 0.0000
f(x*) = 2.6667
Numero de avaliacoes da funcao objetivo: 253
Numero de iteracoes: 11
 		\end{lstlisting}
 		
 		\begin{table}[H]
 			\centering
 			\caption{Entradas atribuídas à função \mcode{varMet} para solução do problema de otimização proposto.}
 			\label{entradasVarMet}
 			\begin{tabular}{|c|c|}
 				\hline
 				\mcode{x0}    & $ [\begin{array}{ccc} 0 & 0 & 0 \end{array}] $ \\ \hline
 				\mcode{theta} & $1$       \\ \hline
 				\mcode{rp0}   & $1$       \\ \hline
 				\mcode{rpInc} & $10$      \\ \hline
 				\mcode{tol}   & $10^{-5}$   \\ \hline
 				\mcode{h}     & $10^{-10}$  \\ \hline
 			\end{tabular}
 		\end{table}
 		
 		Os códigos empregados na obtenção da solução apresentada são introduzidos a seguir e podem ser acessados no link \url{https://github.com/ArthurIasbeck/OTMC4}. 
 		
 		\begin{enumerate}
 			\item Definição da função objetivo (\mcode{funcFile.m}).
 			\begin{lstlisting}
function fObj = funcFile(x)

fObj = 1/3*(x(1) + 1)^3 + x(2);
 			\end{lstlisting}
 				
 			\item Definição do gradiente da função objetivo (\mcode{gradFile.m}).
 			\begin{lstlisting}
function df = gradFile(x)

% O gradiente da funcao objetivo, quando conhecido, deve ser inserido aqui. Caso o mesmo nao possa ser determinado, basta inserir 'df = []'.
df = [];
 			\end{lstlisting}
 			
 			\item Definição das restrições do problemas (\mcode{constFile.m}).
 			\begin{lstlisting}
function [g, h] = constFile(x)

% Cada linha da matriz 'g' representa uma restricao de desigualdade
g = [-x(1) + 1;
     -x(2)];

% Cada linha da matriz 'h' representa uma restricao de igualdade
h = [];
 			\end{lstlisting}
 			
 			\item Código principal (\mcode{main.m})
 			\begin{lstlisting}
% Configuracoes previas 
% Antes de executar a funcao principal 'main.m', eh preciso definir a funcao objetivo, seu gradiente (caso haja), e as restricoes (caso haja). Para tanto basta editar os arquivos 'funcFile.m', 'gradFile.m', e 'constFile.m'

clc; clear; close all;

% Resolvendo o problema de otimizacao
[xOpt, fOpt, nVal, k, alfaValues] = varMet();

% Apresentando os resultados 
for i = 1:length(xOpt)
	fprintf(['x',num2str(i),'* = %.4f\n'], xOpt(i))
end

fprintf('f(x*) = %.4f\n', fOpt)
fprintf('Numero de avaliacoes da funcao objetivo: %d\n', nVal)
fprintf('Numero de iteracoes: %d\n', k)

% Notas 
% A seguir eh introduzida a forma completa da funcao 'varMet'
% [xOpt, fOpt, nVal, k, alfaValues] = varMet(x0, theta, rp0, rpInc, tol, h)

% Para os parametros nao informados pelo usuario serao adotados os valores padrao apresentados a seguir
% x0 = [0 0 0 ... 0]
% theta = 1
% rp0 = 1
% rpInc = 10
% tol = 1e-5
% h = 1e-10
 			\end{lstlisting}
 			
 			\item Implementação do Método das Variáveis Métricas e das demais função utilizadas na computação do mesmo (\mcode{varMet.m}).
 			\begin{lstlisting}
% Implementacao do Metodo da Variavel Metrica
function [xOpt, fOpt, nVal, k, alfaValues] = varMet(x0, theta, rp0, ...
rpInc, tol, h)

% Verificacao do palpite inicial fornecido pelo usuario
n = getOrder;
if nargin < 1 || isempty(x0)
	x0 = zeros(n,1);
else
	% Verificando a consistencia entre x0 e a funcao objetivo
	if n ~= length(x0)
		error('O tamanho de x0 deve ser igual a ordem do problema.');
	end
end

% Verificando entradas da funcao e definindo valores padrao
if nargin < 2 || isempty(theta)
	theta = 1;
end
if nargin < 3 || isempty(rp0)
	rp0 = 1;
end
if nargin < 4 || isempty(rpInc)
	rpInc = 10;
end
if nargin < 5 || isempty(tol)
	tol = 1e-5;
end
if nargin < 6 || isempty(h)
	h = 1e-10;
end

% Verificando se o gradiente analitico foi definido
f = @(x) fObjConst(x);
dfTest = gradFile(x0);
if isempty(dfTest)
	% Gradiente numerico
	if nargin < 4 || isempty(h)
		df = @(x) grad(f,x);
	else
		df = @(x) grad(f,x,h);
	end
	numGrad = 1;
else
	% Gradiente analitico
	df = @(x) gradFile(x);
	numGrad = 0;
end

% Definicao das entradas padrao
if nargin < 2 || isempty(tol)
	tol = 1e-5;
end

if nargin < 3 || isempty(theta)
	theta = 1;
end

% Variaveis para controle de execucao
n = length(x0);
alfaValues = zeros(1,1);
k = 1;
nVal = 0;
H = eye(n);
global rp
rp = rp0;

% Implementacao do processo de otimizacao
while 1
	% Reduzir a dimensao do problema de otimizacao
	g = @(alfa) f(x0 - alfa*H*df(x0));
	
	% Resolver o problema de otimizacao uni-dimensional
	[alfaOpt, gOpt, nVal1] = aureaSec(g,-1,1,1e-4);
	
	% Atualizar a solucao otima
	x = x0 - alfaOpt*H*df(x0);
	
	% Armazenar dados de execucao
	alfaValues(k) = alfaOpt;
	if numGrad
		nVal = nVal + nVal1 + 1;
	else
		nVal = nVal + nVal1 + 2*n;
	end
	
	% OBS : Lembre-se que eh necessaria a computacao do gradiente para atualizacao de x. No entanto, se estivermos utilizando a aproximacao numerica para o gradiente, a computacao do mesmo levara, neste caso a mais avaliacoes da funcao objetivo.
	
	% Verificar a condicao de parada
	cp = norm(x - x0);
	if cp < tol
		break;
	end
	
	% Atualizacao de H (aproximacao para a inversa da Matriz Hessiana)
	p = x - x0;
	y = df(x) - df(x0);
	sigma = p'*y;
	tal = y'*H*y;
	D = ((sigma + theta*tal)/sigma^2)*(p*p') ...
	+ ((theta - 1)/tal)*(H*y)*(H*y)' ...
	- (theta/sigma)*(H*y*p' + p*(H*y)');
	
	H = H + D;
	
	% Atualizar variaveis para a proxima iteracao
	x0 = x;
	k = k + 1;
	rp = rp*rpInc;
end

xOpt = x;
fOpt = f(xOpt);

% Obtencao numerica do gradiente
function df = grad(f, x, h)

% Definindo valores padrao
if nargin < 3
	h = 1e-10;
end

n = length(x);

df = zeros(n,1);
for i = 1:n
	dx = x;
	dx(i) = dx(i) + h;
	df(i) = (f(dx) - f(x))/h;
end

% Implementacao do Metodo da Secao Aurea
function [xOpt, fOpt, k] = aureaSec(f,a,b,tol)
	tal = 0.618;
	
	if nargin < 4
		tol = 1e-8;
	end
	
	alfa = a + (1 - tal)*(b - a);
	beta = a + tal*(b - a);
	fAlfa = f(alfa);
	fBeta = f(beta);
	
	k = 1;
	
	while abs(a-b) > tol
		if fBeta < fAlfa
			a = alfa;
			alfa = beta;
			fAlfa = fBeta;
			beta = a + tal*(b - a);
			fBeta = f(beta);
		elseif fAlfa <= fBeta
			b = beta;
			beta = alfa;
			fBeta = fAlfa;
			alfa = a + (1 - tal)*(b - a);
			fAlfa = f(alfa);
		end
		k = k + 1;
	end
	
	xOpt = (alfa+beta)/2;
	fOpt = f(xOpt);

% Funcao para calculo das penalidades
function P = penalty(x)

% Determinacao dos valores atribuidos a cada uma das restricoes
[g, h] = constFile(x);

% Determinacao das penalidades
Pg = sum(max(0,g).^2);
Ph = sum(h.^2);

% Calculo das penalidades
P = Pg + Ph;

% Funcao objetivo penalizada
function F = fObjConst(x)

global rp

% Computacao da funcao objetivo original
fObj = funcFile(x);

% Computacao das penalidades
P = penalty(x);

% Computacao da funcao objetivo penalizada
F = fObj + rp*P;

% Funcao para verificacao da ordem do problema com base na funcao objetivo
function n = getOrder

x = zeros(1,100);
f0 = funcFile(x);
for i = 1:100
	x(i) = 1;
	f = funcFile(x);
	if f0 == f
		break
	end
	x(i) = 0;
end
n = i - 1;
 			\end{lstlisting}
 		\end{enumerate}
\end{document}